\section{Introduction: the battle of architectures and the four-layer tax}
\label{sec:intro}

\subsection{The von Neumann settlement, and the bottleneck it never removed}

The stored-program computer described in von Neumann's 1945 EDVAC draft
placed instructions and data in one memory addressed by one channel. Its
victory was not a claim of optimality but of \emph{generality}: one memory,
one sequential thread, any program. Every general-purpose machine since is a
descendant. But generality was bought at a structural cost that John Backus
named in his 1977 Turing lecture --- the \emph{von Neumann
bottleneck}~\cite{backus1977}: the processor--memory channel through which
the entire computation must be pumped ``a word at a time.'' Seventy years of
architecture has been, in large part, a campaign against that channel:
caches, prefetch, out-of-order issue, wide vectors, and deep memory
hierarchies are all ways of hiding it, not removing it. The reason it cannot
be removed is now the dominant fact of the field. After the end of Dennard
scaling (c.\ 2006) the binding constraint is energy, and energy in a
von~Neumann machine is dominated not by arithmetic but by \emph{moving
data}: a floating-point operation costs on the order of a picojoule, while
fetching its operands from off-chip memory costs one to two orders of
magnitude more. We compute cheaply and move data dearly, and the
architecture is built around the wrong one.

\subsection{Four subsystems, designed apart}

Underneath the bottleneck lies a deeper and less-discussed fact. A modern
computing system --- and emphatically a \emph{learning} system --- is built
from four subsystems that are conceived, optimised, and verified
\emph{independently}, and then made to coexist:

\begin{enumerate}[leftmargin=1.7em,itemsep=1pt]
  \item the \textbf{interconnect} --- the network-on-chip or fabric topology
        that carries state between compute sites;
  \item the \textbf{error-correcting code} (ECC) --- a separate redundancy
        layer bolted onto memory and links to survive faults;
  \item the \textbf{instruction set} (ISA) --- the algebra of primitive
        operations the machine can perform;
  \item the \textbf{adaptation rule} --- for a learning machine, the
        training/optimisation procedure that changes its parameters.
\end{enumerate}

Each has its own literature, its own optima, and its own designers. Because
they are designed apart, their interfaces must be negotiated, and every
negotiation is a \emph{mismatch tax}, paid three times over:

\begin{table}[H]
\centering
\small
\caption{The four design layers are separate subsystems in every existing
machine, each carrying an independent cost. \TALOS{} pays each cost once, in
one structure.}
\label{tab:four-layers}
\begin{tabular}{@{}p{2.4cm}p{4.6cm}p{5.2cm}@{}}
\toprule
\textbf{Layer} & \textbf{How it is done today} & \textbf{The tax it levies} \\
\midrule
Interconnect & Crossbar / mesh / NoC, designed for bandwidth & Area, static power, routing complexity; all-to-all scales as $\binom{N}{2}$ \\
ECC & Hamming/BCH/LDPC bolted onto storage \& links; surface codes for qubits & Dedicated check bits (+12.5\% for ECC-DRAM) or $\sim\!10^3\times$ physical/logical qubits \\
ISA & Fixed opcode set, co-evolved with a compiler & Semantic gap; the ``von Neumann style'' Backus indicted \\
Adaptation & A separate training loop (back-propagation) over a frozen ISA & Memory for activations/gradients; the train/infer divide; energy $\propto$ data movement \\
\bottomrule
\end{tabular}
\end{table}

\subsection{The challengers, and what each of them fixed}

The alternative architectures were each an attack on \emph{one} of these
four layers, never all four:

\begin{itemize}[leftmargin=1.5em,itemsep=1pt]
  \item \textbf{Harvard} split instruction and data memory --- a
        micro-optimisation of the ISA/interconnect boundary that survives in
        caches and DSPs.
  \item \textbf{Dataflow}~\cite{dennis1975} removed the program counter:
        operations fire when operands are ready. It reformed the
        \emph{instruction} layer (data-driven, not control-driven) but
        offered no principled topology, code, or learning rule.
  \item \textbf{Systolic arrays}~\cite{kung1982} imposed a rhythmic,
        local, nearest-neighbour \emph{interconnect} --- and won the matrix
        engines of today's GPUs and TPUs --- but the topology was an
        engineering choice fitted to one kernel, not a derived structure.
  \item \textbf{Neuromorphic} machines~\cite{mead1990} (TrueNorth, Loihi,
        SpiNNaker) co-located memory and compute and made the substrate
        event-driven and dissipative --- reforming interconnect
        \emph{and} adaptation at once --- but with \emph{ad hoc} topology,
        no native error correction, and no principled state object.
  \item \textbf{Processing-in-memory} and memristive crossbars performed the
        multiply--accumulate \emph{in} the memory array, attacking the
        bottleneck directly, but carried neither dynamics nor a self-model.
  \item \textbf{Quantum} machines made the instruction layer unitary
        (reversible, hence Landauer-cheap) --- and then had to bolt on
        error correction at ruinous overhead, because the code and the
        computation are separate structures that decohere against each other.
\end{itemize}

Every one of these was a genuine advance on one axis. None unified the four,
because none had a \emph{principle} from which a single structure could be
\emph{derived}. The topology was always a guess; the code was always
separate; the learning rule always lived above the ISA. The field has been
searching for a structure that is simultaneously a good network, a good
code, a good instruction algebra, and a good learning rule --- and treating
that as four searches.

\subsection{The thesis: architectural monism}

This document reports that the four searches have one answer, and that the
answer is \emph{forced}, not chosen. Unitary Holonomic Monism (UHM) --- a
mathematical theory that derives the structure of a viable self-modelling
system from a single primitive, the coherence matrix
$\Gam\in\Dens(\mathbb{C}^7)$~\cite{uhm2026} --- proves
that on a seven-mode substrate the interconnect, the error-correcting code,
the instruction algebra, and the learning rule are \emph{not four objects
but one}: the Fano plane $\Fano$. We call this \emph{architectural monism}.
Section~\ref{sec:collapse} states it as two theorems:

\begin{itemize}[leftmargin=1.5em,itemsep=1pt]
  \item \textbf{Theorem~\ref{thm:collapse} (Collapse).} The seven lines of
        $\Fano$ are simultaneously the interconnect; the same incidence is
        the parity-check geometry of the Hamming code $[7,4,3]$; the
        octonionic structure constants supported on those lines are the
        multiplication of the instruction algebra; and the $G_2$-covariant
        fixed point of that algebra is the learning rule. All four are
        functions of one $7\times7$ incidence structure. The mismatch tax is
        identically zero because there are no interfaces to negotiate.
  \item \textbf{Theorem~\ref{thm:unique} (Uniqueness).} This collapse is
        realisable at a \emph{unique} node count, $N=7$. The coincidence
        ``projective-plane order $=$ Hamming length'' recurs (at $N = 7,
        31, \dots$), but only at $N=7$ does the instruction algebra also
        carry a non-trivial associator --- the third-order structure on
        which the whole design rests. Non-associativity breaks the tie in
        favour of seven, and nowhere else.
\end{itemize}

The architecture that follows, \TALOS{},\footnote{After the bronze
automaton of Greek myth, the first artificial being --- here the physical
body of an artificial mind. The name is a codename, not an acronym.} is
therefore not a proposal among proposals. It is the hardware form of a
mathematical object. Its state is a $784$-byte density matrix that fits in
L1 cache; its one instruction is the UHM Lindbladian $\LindOmega$; its
network is seven buses; its fault tolerance is the network; its learning is
the third term of the instruction; and its power draw tracks how fast it is
learning, not how hard it is computing.

\subsection{Honest boundaries, stated first}

Two theorems of UHM constrain what \TALOS{} can and cannot claim, and we
foreground both so that nothing below is read as more than it is.

\begin{itemize}[leftmargin=1.5em,itemsep=1pt]
  \item \textbf{Substrate multiplicity (T-153).} UHM's substrate-closure
        theorem proves the computation is substrate-\emph{independent}: the
        choice of physical embedding is a pure $G_2$ gauge. There is no
        unique physical machine --- only a unique \emph{logical}
        architecture and a family of substrates that realise it. \TALOS{}
        is the impedance-matched member of that family
        (\S\ref{sec:realizability}), not the sole possibility. This is a
        strength (portability), not a weakness.
  \item \textbf{Computational bound (BQP).} UHM proves the regeneration term
        $\Regen$ grants \emph{no} speed-up beyond the complexity class BQP.
        \TALOS{} is not a faster-than-quantum machine. Its advantage is in
        energy, adaptivity, and fault-tolerance --- capability per watt and
        native learning --- never in complexity class. We repeat this wherever
        a reader might over-read a claim.
\end{itemize}

The remainder of the document is a specification. Section~\ref{sec:collapse}
proves the two theorems. Sections~\ref{sec:substrate}--\ref{sec:composition}
specify the state, the fabric, the code-as-topology, the energy model, the
self-model, and the composition law. Section~\ref{sec:realizability} gives a
realisability ladder from today's GPU emulation to a driven-dissipative
analogue device. Section~\ref{sec:abi} specifies the application binary
interface to \SYNARC{}, the software mind for which \TALOS{} is the body.
Section~\ref{sec:redteam} subjects the design to an explicit red-team, and
Section~\ref{sec:benchmarks} reduces it to falsifiable engineering
measurements against a von~Neumann baseline. Every number is machine-checked.
